\documentclass[10pt]{article}

\usepackage[margin=0.75in]{geometry}
\usepackage{amsmath, amssymb}
\usepackage{booktabs}
\usepackage{hyperref}
\usepackage{enumitem}
\setlist{nosep,leftmargin=*}
\setlength{\parskip}{0.25em}
\setlength{\parindent}{0pt}

\title{\textbf{Project Proposal}\\
Scalable Graph Partitioning for Communication-Efficient Distributed Optimization}
\author{Team Members: Name 1, Name 2, Name 3}
\date{CS240: Algorithm Design and Analysis, Spring 2026}

\begin{document}

\maketitle

\section{Topic Selection}

We choose \textbf{Topic 2: Community Detection / Graph Partitioning} from the
reference list, but use a different modern application:
\textbf{communication-efficient distributed optimization}. In distributed
optimization, a large-scale problem is solved by multiple agents or workers
that exchange information through a communication graph. The graph structure
directly affects communication cost, load balance, and convergence behavior.
Instead of a standard social-network community detection setting, we study how
classical graph partitioning algorithms can organize large-scale distributed
computation. This application is closely related to large-scale optimization,
multi-agent systems, and distributed machine learning.

\section{Problem Formulation}

Let $G=(V,E,w)$ be an undirected weighted graph, where each vertex represents a
data block, variable block, agent, or computational task, and each edge
represents communication or dependency between two vertices. Given an integer
$k$, the goal is to partition $V$ into $k$ disjoint subsets
$V_1,\ldots,V_k$.

The optimization objective is
\[
    \min_{V_1,\ldots,V_k}
    \sum_{(u,v)\in E,\; u\in V_i,\; v\in V_j,\; i\neq j} w_{uv},
\]
subject to a balance constraint
\[
    |V_i| \leq (1+\epsilon)\frac{|V|}{k},
    \quad i=1,\ldots,k.
\]

Here the cut objective models cross-partition communication cost, and the
balance constraint models computational load balance. The input is a weighted
graph and the desired number of partitions; the output is a balanced
$k$-partition of the vertices.

\section{Related Work}

Graph partitioning is a classical algorithmic problem in scientific computing,
parallel computation, and large-scale graph processing. Spectral partitioning
uses eigenvectors of the graph Laplacian to find balanced cuts. The
Kernighan--Lin algorithm is a classical local-search method that iteratively
swaps vertices to reduce cut size. Multilevel methods such as METIS improve
scalability by graph coarsening, partitioning, and refinement. Distributed
optimization over networks is closely related: in consensus and distributed
gradient methods, agents update local variables while communicating with
neighbors, so topology affects both convergence and communication overhead.

We plan to use the \textbf{Kernighan--Lin graph partitioning heuristic} as the
specific paper/method selected for reproduction. Spectral bisection and a
greedy balanced partitioning method will be used as baselines or variants.
Related papers include:
\begin{itemize}
    \item B. W. Kernighan and S. Lin, ``An Efficient Heuristic Procedure for
    Partitioning Graphs,'' \emph{Bell System Technical Journal}, 1970.
    DOI: \href{https://doi.org/10.1002/j.1538-7305.1970.tb01770.x}
    {10.1002/j.1538-7305.1970.tb01770.x}.
    \item G. Karypis and V. Kumar, ``A Fast and High Quality Multilevel Scheme
    for Partitioning Irregular Graphs,'' \emph{SIAM Journal on Scientific
    Computing}, 1998. DOI: \href{https://doi.org/10.1137/S1064827595287997}
    {10.1137/S1064827595287997}.
    \item A. Nedic and A. Ozdaglar, ``Distributed Subgradient Methods for
    Multi-Agent Optimization,'' \emph{IEEE Transactions on Automatic Control},
    2009. DOI: \href{https://doi.org/10.1109/TAC.2008.2009515}
    {10.1109/TAC.2008.2009515}.
\end{itemize}

\section{Algorithm and Technical Plan}

We plan to implement and compare the following graph partitioning methods:
\begin{enumerate}
    \item \textbf{Greedy balanced partitioning baseline.} Vertices are assigned
    to partitions one by one while maintaining balance and reducing local cut
    cost.
    \item \textbf{Spectral bisection.} We compute the graph Laplacian, use the
    Fiedler vector to bisect the graph, and recursively apply bisection to
    obtain $k$ partitions.
    \item \textbf{Kernighan--Lin refinement.} Starting from an initial
    partition, we iteratively swap vertices across partitions to reduce the cut
    value while respecting balance constraints.
\end{enumerate}

The planned implementation will be in Python, using standard scientific
computing tools such as NumPy, SciPy, and NetworkX. The planned evaluation
setup will include grid graphs, Erdos--Renyi random graphs, random geometric
graphs, and possibly one small real network. We plan to analyze each method's
time complexity and compare cut size, balance ratio, runtime, and scalability
as $|V|$ and $|E|$ increase. As an optional application-level experiment, we
plan to simulate consensus averaging or distributed gradient descent and
measure objective gap, consensus error, and communication cost under different
partitions.

\section{Project Scope and Expected Goals}

The minimum goal is to reproduce the Kernighan--Lin heuristic, implement
spectral bisection and a greedy baseline, and compare their partition quality
and runtime. This directly demonstrates how classical graph algorithms can
address a modern distributed optimization bottleneck: communication-efficient
work allocation.

If time permits, we will add one or more extensions:
\begin{itemize}
    \item weighted graph partitioning with nonuniform communication costs;
    \item comparison with an existing library such as NetworkX or METIS;
    \item a stronger distributed optimization simulation with different
    network topologies and communication budgets.
\end{itemize}

\section{Team Information}

\begin{tabular}{@{}lll@{}}
\toprule
Member & Responsibility & Expected Contribution \\
\midrule
Name 1 & Algorithm implementation & Spectral partitioning and experiments \\
Name 2 & Theory and writing & Problem formulation and complexity analysis \\
Name 3 & Evaluation & Distributed optimization simulation and plots \\
\bottomrule
\end{tabular}

\end{document}
